\documentclass[11pt]{article}

\usepackage[letterpaper,left=1.25in,right=1.25in,top=1in,bottom=1in]{geometry}
\usepackage{setspace}
\usepackage{microtype}
\usepackage{amsmath,amssymb,amsfonts,amsthm}
\usepackage{mathtools}
\usepackage{hyperref}
\usepackage{graphicx}
\usepackage{booktabs}
\emergencystretch=1em
\sloppy

\usepackage{fontspec}
\usepackage{unicode-math}

\setmainfont{Latin Modern Roman}


% Dispatcher: \MSOglyph{Ocularo} expands to \msoOcularo, etc.
\newcommand{\MSOglyph}[1]{\csname mso#1\endcsname}

% MSO symbol assignments (pure math glyphs)
\newcommand{\msoOcularo}{\ensuremath{\odot}}
\newcommand{\msoMonoturno}{\ensuremath{\circlearrowleft}}
\newcommand{\msoLamphron}{\ensuremath{\rightsquigarrow}}
\newcommand{\msoLamphrodyno}{\ensuremath{\rightleftharpoons}}
\newcommand{\msoGravito}{\ensuremath{\Delta}}
\newcommand{\msoQuarko}{\ensuremath{\varepsilon}}
\newcommand{\msoElectro}{\ensuremath{\mathrm{e}}}
\newcommand{\msoLepto}{\ensuremath{\ell}}
\newcommand{\msoFermio}{\ensuremath{\mathrm{f}}}
\newcommand{\msoPhoto}{\ensuremath{\gamma}}
\newcommand{\msoGluo}{\ensuremath{\mathrm{g}}}
\newcommand{\msoMeso}{\ensuremath{\mu}}
\newcommand{\msoBoso}{\ensuremath{\beta}}
\newcommand{\msoProto}{\ensuremath{\mathrm{p}}}
\newcommand{\msoNeutro}{\ensuremath{\circ}}
\newcommand{\msoBaryo}{\ensuremath{\mathrm{B}}}
\newcommand{\msoHadro}{\ensuremath{\mathrm{H}}}
\newcommand{\msoAto}{\ensuremath{\alpha}}
\newcommand{\msoMoleculo}{\ensuremath{\mathcal{M}}}
\newcommand{\msoUnicellulo}{\ensuremath{\mathcal{U}}}
\newcommand{\msoCrystolattisso}{\ensuremath{\Lambda}}
\newcommand{\msoMicro}{\ensuremath{\sigma}}
\newcommand{\msoMegalobio}{\ensuremath{\mathcal{B}}}
\newcommand{\msoGaiama}{\ensuremath{\mathcal{E}}}
\newcommand{\msoTechnio}{\ensuremath{\mathcal{T}}}
\newcommand{\msoTella}{\ensuremath{\oplus}}
\newcommand{\msoGaiaselena}{\ensuremath{\circledcirc}}
\newcommand{\msoSolara}{\ensuremath{\partial}}
\newcommand{\msoPerinebula}{\ensuremath{\star}}
\newcommand{\msoVacuo}{\ensuremath{\varnothing}}
\newcommand{\msoOriocygnobrachio}{\ensuremath{\Omega}}
\newcommand{\msoGalacto}{\ensuremath{\mathcal{G}}}
\newcommand{\msoProximasystada}{\ensuremath{\circledast}}
\newcommand{\msoVirgo}{\ensuremath{\mathrm{V}}}
\newcommand{\msoLaniakeasuperclusto}{\ensuremath{\bowtie}}
\newcommand{\msoPiscescetusfilamento}{\ensuremath{\longrightarrow}}
\newcommand{\msoCytosis}{\ensuremath{\nabla}}


\setstretch{1.12}
\setlength{\parindent}{0pt}
\setlength{\parskip}{0.75em}

\newtheorem{definition}{Definition}
\newtheorem{proposition}{Proposition}
\newtheorem{remark}{Remark}

\title{The Mereological Space Ontology:\\
Glyphs, Scales, and the Observer in a Layered Part--Whole Universe}
\author{Flyxion}
\date{\today}

\begin{document}
\maketitle

\begin{abstract}
This essay develops a compact formal and interpretive account of a ``Mereological Space Ontology'' (MSO) in which the universe is presented as an ordered family of part--whole decompositions across scales, indexed by a reversible glyph-string. The MSO is treated as a two-sided construction: an ontic ordering of compositional strata and an epistemic ordering of observational access. The glyphs are interpreted as operators that mark regime-changes in granularity, binding, and instrumentation, rather than as mere pictograms. A central aim is to clarify how the same chain can be read forward as a refinement of composition and backward as a narrative of discovery, while remaining consistent with a single underlying mereological scaffold.

Formally, the MSO is framed as a scale-indexed family of mereological structures equipped with coarse-graining maps and boundary conditions that encode when adjacent regimes require a change of primitives. Stability of wholes is treated as regime-relative, permitting both dynamical and representational notions of persistence, and motivating a distinction between ontic parthood and instrument-induced partition. The essay develops several mutually reinforcing lenses---measure-theoretic persistence, spectral and topological diagnostics of regime transition, and categorical and algebraic reformulations---to make precise the sense in which ``wholes'' can remain coherent while their admissible descriptions change. The resulting picture is neither strict reductionism nor unconstrained holism, but a disciplined account of layered composition under constraint in which scale, stability, and observation jointly determine what counts as a unit.
\end{abstract}

\newpage
\section{Motivation and Setup}

The MSO begins from a simple premise: any world-description that treats ``things'' as stable must implicitly choose a decomposition of the world into parts, and must silently decide which compositions count as units. Mereology provides a language for this choice. The MSO adds two ingredients. First, it treats scale not as an afterthought but as a primary index that organizes admissible parthood relations. Second, it encodes salient regime transitions through a glyph-string that can be read both as an ontological ladder of composition and as an epistemic ladder of exploration.

The long composite token
\begin{flushleft}

\textsc{ocularomonoturnolamphrolamphrodynogravitoquarkoelectroleptofermiophoto%
\allowbreak gluomesobosoprotoneutrobaryohadroatomoleculounicellulocrystolattisso%
\allowbreak micromegalobiogaiamatechniotellagaiaselenasolaraperinebulavacuo%
\allowbreak oriocygnobrachiogalactoproximasystadavirgolaniakeasuperclusto%
\allowbreak piscescetusfilamentocytosis}

\end{flushleft}

functions as a compression of the scale spectrum, from observer-marked space through microphysical carriers, chemistry, biology, planet-scale organization, technological accretion, stellar and galactic structure, and finally cosmic web and horizon-structure. 

\begin{figure}[htbp]
\centering
\[
\MSOglyph{Ocularo}%
\MSOglyph{Monoturno}%
\MSOglyph{Lamphron}%
\MSOglyph{Lamphrodyno}%
\MSOglyph{Gravito}%
\MSOglyph{Quarko}%
\MSOglyph{Electro}%
\MSOglyph{Lepto}%
\MSOglyph{Fermio}%
\MSOglyph{Photo}%
\MSOglyph{Gluo}%
\MSOglyph{Meso}%
\MSOglyph{Boso}%
\MSOglyph{Proto}%
\MSOglyph{Neutro}%
\MSOglyph{Baryo}%
\MSOglyph{Hadro}%
\MSOglyph{Ato}%
\MSOglyph{Moleculo}%
\MSOglyph{Unicellulo}%
\MSOglyph{Crystolattisso}%
\MSOglyph{Micro}%
\MSOglyph{Megalobio}%
\MSOglyph{Gaiama}%
\MSOglyph{Technio}%
\MSOglyph{Tella}%
\MSOglyph{Gaiaselena}%
\MSOglyph{Solara}%
\MSOglyph{Perinebula}%
\MSOglyph{Vacuo}%
\MSOglyph{Oriocygnobrachio}%
\MSOglyph{Galacto}%
\MSOglyph{Proximasystada}%
\MSOglyph{Virgo}%
\MSOglyph{Laniakeasuperclusto}%
\MSOglyph{Piscescetusfilamento}%
\MSOglyph{Cytosis}
\]
\caption{MSO glyph sequence (ontic ascent) from observer-indexed partition to ontological threshold.}
\label{fig:mso-forward}
\end{figure}

\section{Mnemonic and Didactic Function of the Composite Token}

Before turning to formal structure, it is useful to clarify the pedagogical intention of the composite token that encodes the MSO scale spectrum. Its construction is deliberately mnemonic. Much as certain extended words in poetic or musical traditions serve to compress complex ideas into rhythmically repeatable form, the MSO token is designed to be pronounceable in a single breath and segmentable by cadence. Its length is not ornamental excess but instructional device.

The inspiration is partly phonetic and partly cognitive. A sufficiently repeated composite utterance can function as a scaffold for thought: each internal segment cues a regime, and each regime cues a distinct family of invariants and admissible wholes. By articulating the token repeatedly, one rehearses the ordered passage through hierarchical scales, from microphysical mediation through chemical, biological, planetary, technological, stellar, and cosmological organization. The act of recitation thus becomes a traversal of the scale ladder.

Because the token can be segmented rhythmically, it permits counting off its internal strata on one’s fingers. This bodily indexing reinforces the ordered structure and gradually renders the sequence automatic. Over time, the utterance ceases to feel like an arbitrary string and instead becomes a compressed map of the layered ontology. In this way, the composite word serves not merely as a label but as a didactic engine: it primes holistic thinking by activating universal constructs across hierarchical scales, integrating them into a single rehearsable arc.

\section{Mereological Core}

Let $U$ denote a universe of discourse and let $\leq$ be a parthood relation on $U$. In the minimal setting, $\leq$ is taken to be reflexive and transitive; when antisymmetry is assumed, $(U,\leq)$ is a poset. Write $x \sqcup y$ for a fusion (sum) operation when it exists, and write $x \perp y$ to indicate disjointness in the sense of no shared part.

\begin{definition}[Scale-indexed mereology]
A \emph{scale-indexed mereology} is a family $\{(U_s,\leq_s)\}_{s\in S}$ where $S$ is a totally ordered set of ``scales,'' and where each $(U_s,\leq_s)$ is a parthood structure at scale $s$. For $s \leq t$ in $S$, a \emph{coarsening map} $\pi_{t\to s}:U_t\to U_s$ identifies fine-scale entities as parts of coarse-scale units.
\end{definition}

The MSO does not require that each $U_s$ contains the same objects. It instead treats each $U_s$ as a regime-appropriate ontology with its own admissible units. This is crucial because the transition from, for example, quark-level descriptions to hadronic descriptions is not merely a change in magnification; it is a change in what counts as stable.

\begin{definition}[Regime boundary]
A \emph{regime boundary} is a scale transition $s\to t$ such that no single parthood language is simultaneously stable across both scales, so that the coarsening map $\pi_{t\to s}$ must be accompanied by a change of primitives, constraints, or effective laws.
\end{definition}

The glyphs in the MSO are best read as markers of such regime boundaries and their corresponding stabilizations.

\section{Glyphs as Operators, Not Decorations}

A glyph in the MSO is not merely a label for ``an object-type.'' It is a compact operator that signals a stable descriptive regime: what the units are, which compositions are treated as primitive, and which instruments mediate access. In this sense the glyph-string is a presentation of a stratified ontology.

The following section gives a detailed reading of each glyph as a regime marker. The descriptions are phrased so that each glyph can be understood as (i) an ontic stabilizer of units at a certain granularity and (ii) an epistemic cue about how that granularity is accessed.

\subsection*{Glyph Interpretations}

\paragraph{$\MSOglyph{Ocularo}$ ; Observer-marked phenomenal space.}
This glyph denotes the primordial observational cut. Any articulation of wholes begins from a selection of contrasts, variables, and invariants that count as given. It marks the epistemic root of the MSO: the universe as partitioned through attention, resolution, and instrumentally mediated access.

\paragraph{$\MSOglyph{Monoturno}$ ; Scale traversal and structured refinement.}
This glyph signifies ordered passage between regimes. Movement across scales is not mere magnification but the application of projection operators and consistency-preserving maps. It encodes the functorial relation between adjacent ontologies.

\paragraph{$\MSOglyph{Lamphron}$ ; Radiant mediation.}
This glyph denotes regimes in which access and structure are revealed through propagating carriers. Illumination, emission, and radiative inference become primary epistemic tools.

\paragraph{$\MSOglyph{Lamphrodyno}$ ; Dynamical exchange and coupling.}
This glyph marks interaction as constitutive. Stable configurations arise from mediated exchange processes that coordinate parts into coherent regimes.

\paragraph{$\MSOglyph{Gravito}$ ; Constrained motion and field response.}
This glyph denotes matter insofar as it manifests through response to forces. Attraction, repulsion, confinement, and resonance define stable trajectories and binding relations.

\paragraph{$\MSOglyph{Quarko}$ ; Substructural resolution.}
This glyph indicates regimes revealed through fine-grained instrumentation. What counts as a unit shifts with enhanced resolving power, altering the admissible invariants.

\paragraph{$\MSOglyph{Electro}$ ; Charged persistence.}
This glyph marks the regime in which stored potential and charge separation become structuring principles. It emphasizes persistence through energetic differentiation.

\paragraph{$\MSOglyph{Lepto}$ ; Lightness and minimal quanta.}
This glyph designates the lightest stabilized excitations in a regime. It marks thresholds where mass, charge, or interaction strength approaches minimality.

\paragraph{$\MSOglyph{Fermio}$ ; Fermionic individuation.}
This glyph denotes discrete matter constituents whose exclusion and persistence generate compositional language.

\paragraph{$\MSOglyph{Photo}$ ; Photonic revelation.}
This glyph emphasizes luminous quanta as carriers of information and structure.

\paragraph{$\MSOglyph{Gluo}$ ; Binding remainder and confining structure.}
This glyph marks regimes in which constituents are inseparable without altering identity. It signals effective parameters required for structural coherence.

\paragraph{$\MSOglyph{Meso}$ ; Mesoscopic targeting.}
This glyph denotes intermediate stabilization, where collective behavior becomes more relevant than elementary composition.

\paragraph{$\MSOglyph{Boso}$ ; Bosonic field coherence.}
This glyph represents collective excitation and wave-mediated ordering, where continuity supersedes discreteness.

\paragraph{$\MSOglyph{Proto}$ ; Protonic charge core.}
This glyph marks stable positive cores around which larger composites are organized.

\paragraph{$\MSOglyph{Neutro}$ ; Neutral core stability.}
This glyph denotes balancing constituents that stabilize composites without contributing net charge.

\paragraph{$\MSOglyph{Baryo}$ ; Baryonic massiveness.}
This glyph emphasizes heavy composite matter whose persistence grounds higher-level aggregation.

\paragraph{$\MSOglyph{Hadro}$ ; Hadronic assembly.}
This glyph marks bound composites formed through strong interaction regimes.

\paragraph{$\MSOglyph{Ato}$ ; Atomic machinery.}
This glyph signals the emergence of discrete atomic units capable of stable recombination.

\paragraph{$\MSOglyph{Moleculo}$ ; Molecular transformation.}
This glyph denotes controlled combinatorial chemistry and classification through reaction invariants.

\paragraph{$\MSOglyph{Unicellulo}$ ; Minimal cellular closure.}
This glyph marks the transition to self-maintaining biological organization defined thermodynamically rather than merely structurally.

\paragraph{$\MSOglyph{Crystolattisso}$ ; Crystalline periodicity.}
This glyph denotes lattice symmetry and ordered repetition as stability principles.

\paragraph{$\MSOglyph{Micro}$ ; Microscopic inspection.}
This glyph represents refined observational access within structured domains.

\paragraph{$\MSOglyph{Megalobio}$ ; Large-scale biological integration.}
This glyph denotes macroscopic life forms stabilized by multi-layered constraint integration.

\paragraph{$\MSOglyph{Gaiama}$ ; Biospheric generativity.}
This glyph marks life as a planetary-layer process sustained by cycles and feedback.

\paragraph{$\MSOglyph{Technio}$ ; Engineered closure.}
This glyph denotes artificial systems whose identity persists under designed substitution and interface constraints.

\paragraph{$\MSOglyph{Tella}$ ; Planetary individuation.}
This glyph fixes a planet as a coupled dynamical whole composed of interacting subsystems.

\paragraph{$\MSOglyph{Gaiaselena}$ ; Satellite coupling.}
This glyph denotes orbital partnership and resonance-based stabilization.

\paragraph{$\MSOglyph{Solara}$ ; Stellar energetic hub.}
This glyph marks stars as organizing centers structuring dependent systems through radiative output.

\paragraph{$\MSOglyph{Perinebula}$ ; Nebular precursors.}
This glyph denotes diffuse generative regimes in which future wholes incubate.

\paragraph{$\MSOglyph{Vacuo}$ ; Vacuum substrate.}
This glyph treats emptiness as structured separation enabling propagation and interaction.

\paragraph{$\MSOglyph{Oriocygnobrachio}$ ; Galactic-arm indexing.}
This glyph denotes large-scale structural segmentation within galactic organization.

\paragraph{$\MSOglyph{Galacto}$ ; Galactic aggregation.}
This glyph marks gravitationally bound stellar systems as stable macro-units.

\paragraph{$\MSOglyph{Proximasystada}$ ; System-level architecture.}
This glyph generalizes planetary individuation to multi-body configurations governed by formation history.

\paragraph{$\MSOglyph{Virgo}$ ; Cluster partition.}
This glyph denotes large-scale grouping indexed through cosmographic classification.

\paragraph{$\MSOglyph{Laniakeasuperclusto}$ ; Supercluster connectivity.}
This glyph marks web-like large-scale structure where connectivity defines ontological coherence.

\paragraph{$\MSOglyph{Piscescetusfilamento}$ ; Filamentary continuity.}
This glyph foregrounds long-range structural strands that bind vast regions into integrated wholes.

\paragraph{$\MSOglyph{Cytosis}$ ; Ontological threshold.}
This glyph marks the boundary of articulation, the horizon beyond which current invariants fail to stabilize describable units.

\section{Forward and Reverse Readings}

The forward traversal, culminating in $\MSOglyph{Cytosis}$, may be read as ontic ascent: beginning from observer-indexed partition and passing through interaction, composition, biological closure, planetary integration, and finally cosmological connectivity. Each transition introduces a new stability criterion and a new admissible parthood relation.

The reverse traversal, beginning at $\MSOglyph{Cytosis}$, functions epistemically. It narrates discovery from horizon-scale structure inward through cosmic aggregation, stellar and planetary access, technological mediation, biological organization, chemical composition, and ultimately microphysical interaction, returning to $\MSOglyph{Ocularo}$. The reversibility expresses a deeper claim: that the same scale-indexed mereological order supports both the construction of wholes and the acquisition of wholes, provided one distinguishes ontic parthood from instrument-induced partition.

\noindent
\ttfamily
Piscescetusfilamento\allowbreak
laniakea\allowbreak
virgo\allowbreak
superclusto\allowbreak
proximasystada\allowbreak
galacto\allowbreak
oriocygnobrachio\allowbreak
perivacuo\allowbreak
nebula\allowbreak
solara\allowbreak
gaiaselena\allowbreak
tella\allowbreak
technio\allowbreak
gaiama\allowbreak
megalo\allowbreak
microbio\allowbreak
crystolattisso\allowbreak
unicellulo\allowbreak
moleculo\allowbreak
ato\allowbreak
hadro\allowbreak
baryo\allowbreak
neutro\allowbreak
proto\allowbreak
boso\allowbreak
meso\allowbreak
gluo\allowbreak
photo\allowbreak
fermio\allowbreak
lepto\allowbreak
electro\allowbreak
quarko\allowbreak
gravito\allowbreak
lamphrodyno\allowbreak
lamphro\allowbreak
monoturno\allowbreak
ocularo\allowbreak
cytosis

\normalfont

\begin{remark}
In MSO terms, the ``reverse ontology'' is not the negation of the forward ontology. It is the contravariant action of the observational process. If $S$ is ordered coarse-to-fine ontically, then epistemic access often runs fine-to-coarse: instruments first present signals, then stabilize objects, then stabilize large-scale structure.
\end{remark}
\begin{figure}[htbp]
\centering
\[
\MSOglyph{Piscescetusfilamento}%
\MSOglyph{Laniakeasuperclusto}%
\MSOglyph{Virgo}%
\MSOglyph{Proximasystada}%
\MSOglyph{Galacto}%
\MSOglyph{Oriocygnobrachio}%
\MSOglyph{Vacuo}%
\MSOglyph{Perinebula}%
\MSOglyph{Solara}%
\MSOglyph{Gaiaselena}%
\MSOglyph{Tella}%
\MSOglyph{Technio}%
\MSOglyph{Gaiama}%
\MSOglyph{Megalobio}%
\MSOglyph{Micro}%
\MSOglyph{Crystolattisso}%
\MSOglyph{Unicellulo}%
\MSOglyph{Moleculo}%
\MSOglyph{Ato}%
\MSOglyph{Hadro}%
\MSOglyph{Baryo}%
\MSOglyph{Neutro}%
\MSOglyph{Proto}%
\MSOglyph{Boso}%
\MSOglyph{Meso}%
\MSOglyph{Gluo}%
\MSOglyph{Photo}%
\MSOglyph{Fermio}%
\MSOglyph{Lepto}%
\MSOglyph{Electro}%
\MSOglyph{Quarko}%
\MSOglyph{Gravito}%
\MSOglyph{Lamphrodyno}%
\MSOglyph{Lamphron}%
\MSOglyph{Monoturno}%
\MSOglyph{Ocularo}%
\MSOglyph{Cytosis}
\]
\caption{Reverse MSO glyph traversal terminating at the ontological threshold.}
\label{fig:mso-reverse-chain}
\end{figure}


\section{A Minimal Formal Reading of the Glyph-String}

A concise way to formalize the MSO glyph-string is to treat it as a word in an alphabet $\Sigma$ of regime operators acting on a base observational cut. Let $\mathcal{O}$ denote an observational state (a chosen partition of appearances) and let each glyph $g\in\Sigma$ denote an operator $T_g$ that changes the admissible unit language, instruments, or constraints. The forward MSO chain then defines a composite transformation
\[
\mathcal{O} \xrightarrow{T_{\��️}} \mathcal{O}_1 \xrightarrow{T_{\��}} \mathcal{O}_2 \xrightarrow{}\cdots\xrightarrow{T_{\��}} \mathcal{O}_n.
\]
The reverse chain is the composite in the reverse order. These composites need not be inverses in the strict algebraic sense, because regime transitions are typically lossy: coarsening forgets microdetail, while refinement introduces new degrees of freedom. Nevertheless, the MSO asserts that both composites are coherent traversals of the same stratified ontology, and that coherence is maintained by requiring compatibility conditions between adjacent regimes.

\begin{proposition}[Compatibility as a boundary condition]
If each adjacent regime transition $s\to t$ admits a coarsening map $\pi_{t\to s}$ and a constraint set $C_s$ defining admissible wholes at scale $s$, then MSO coherence requires that images of admissible fine-scale wholes satisfy coarse-scale constraints, i.e.,
\[
x\in U_t \text{ admissible} \quad \Longrightarrow \quad \pi_{t\to s}(x)\in U_s \text{ admissible}.
\]
\end{proposition}

This proposition encodes the intuitive demand that the ladder not contradict itself: when one declares a whole at a fine scale, its coarse shadow should not violate the constraints that define wholes at the coarse scale.

\section{From Parts to Wholes}

The MSO glyph-string is best understood as a compact presentation of scale-indexed mereology together with a theory of observational cuts. Each glyph names a stabilization of units and a boundary condition on what counts as a part, what counts as a whole, and what kinds of instruments or invariants justify the carving. The forward reading functions as an ontic refinement of composition up to horizons and filamentary connectivity; the reverse reading functions as an epistemic narrative of discovery from boundary and large-scale structure down to microphysical mediation and finally back to the observer. The philosophical point is that the same universe can legitimately be read in both directions only if one keeps distinct the mereology of the world and the mereology induced by observation, naming, and tool-mediated access.

In the remainder of this essay, the conceptual framework of the Mereological Space Ontology will be recast in progressively more formal terms. What has thus far been presented as a stratified account of parts, wholes, and regime transitions will be interpreted through measure-theoretic stability criteria, spectral and topological diagnostics, categorical structure, and finally an algebraic reformulation. These mathematical interpretations are not intended as ornamental abstractions, but as clarifications of the structural commitments implicit in the MSO. Each formal lens isolates a distinct aspect of regime organization, thereby rendering precise the conditions under which wholes persist, transitions occur, and scale-indexed descriptions remain coherent.

\section{Mereology and Regime Change}

A central claim of the Mereological Space Ontology (MSO) is that ontological stratification is not merely quantitative but qualitative. Scale transitions are not simple magnifications. They are regime changes in admissible part--whole relations.

Let $(U_s,\leq_s)$ and $(U_t,\leq_t)$ be adjacent scales with $s < t$. The MSO asserts that there exist cases in which no embedding functor $F : U_s \to U_t$ preserves all mereological structure. That is, certain wholes at scale $s$ cannot be recovered purely by aggregation of parts at scale $t$ without introducing new organizing constraints.

This motivates a distinction between:

\begin{definition}[Compositional sufficiency]
A scale $t$ is compositionally sufficient for scale $s$ if every whole in $U_s$ is the image under coarsening of some fusion in $U_t$.
\end{definition}

\begin{definition}[Constraint emergence]
A regime transition exhibits constraint emergence if there exist wholes in $U_s$ whose defining invariants are not reducible to parthood relations in $U_t$ alone, but require new structural predicates.
\end{definition}

Chemical bonding relative to atomic structure, biological organization relative to chemistry, and planetary ecology relative to geology are candidate examples. The MSO does not deny microphysical grounding; rather, it denies that grounding alone exhausts the ontology of wholes.

\section{Observer-Indexed Ontology}

The glyph corresponding to the observer signals that any ontology is implicitly observer-indexed. Let $\mathcal{O}$ denote an observational partition of $U$. Then:

\begin{definition}[Observer-indexed parthood]
Given an observational structure $\mathcal{O}$, define $x \leq_{\mathcal{O}} y$ if and only if $x$ is treated as a part of $y$ under the invariants preserved by $\mathcal{O}$.
\end{definition}

This makes explicit that parthood is not merely a primitive relation in the world, but a relation stabilized under a choice of observables. Two observational cuts may induce distinct, though compatible, mereologies.

The MSO therefore distinguishes between ontic mereology and epistemic mereology. Ontic mereology refers to parthood relations taken to obtain independently of any particular act of observation or measurement. In this sense, it concerns the structural organization of the world as described by the best available dynamical or field-theoretic account, and it treats wholes as stabilized configurations of underlying processes whether or not they are directly accessible to a given observer.

Epistemic mereology, by contrast, refers to parthood relations induced by the constraints of observation, instrumentation, and resolution. Here, wholes are defined relative to invariants preserved under a specific observational cut. What counts as a part or a whole depends upon which degrees of freedom are resolved, which quantities are treated as noise, and which transformations are considered admissible within the representational framework.

The forward and reverse glyph strings encode these two orientations. The forward reading may be interpreted as ontic ascent, articulating progressively more inclusive or structurally dominant regimes of organization within the world itself. The reverse reading may be interpreted as epistemic descent, tracing the path by which an observer, equipped with instruments and models, reconstructs larger-scale structure from progressively refined data. Together, they express the bidirectional compatibility between a world-structured mereology and an observer-indexed mereology grounded in representation.

\section{Filamentary Structure and Non-Local Wholes}

At the largest scales, the MSO introduces filamentary connectivity as a primary ontological mode. Traditional mereology often presupposes spatial contiguity as a criterion for parthood. Cosmic filaments challenge this assumption.

\begin{definition}[Non-local whole]
A whole $w$ is non-local if its parts are not contiguous in Euclidean space, but are bound by dynamical or structural constraints that preserve a higher-level invariant.
\end{definition}

Examples include gravitationally bound superclusters, information networks, and technological infrastructures. In such cases, connectivity replaces adjacency as the operative principle of fusion.

This motivates an extension of classical mereology to include relational fusion:

\[
x \sqcup_R y
\]

where $R$ is a binding relation (gravitational, informational, functional). The MSO thereby generalizes parthood beyond simple material adjacency.

\section{Vacuum as Ontological Condition}

The glyph marking vacuum does not denote absence, but structural condition. In MSO terms, vacuum is the regime that permits separability.

Let $V$ denote the vacuum regime. Then for distinct wholes $a,b$ to be treated as disjoint, there must exist a region $v \in V$ such that:

\[
a \perp v \perp b.
\]

Vacuum thus functions as a separator enabling individuated wholes. It is not a thing among things, but a boundary-enabler in the ontology of separation.

\section{Technological Accretion and Artificial Mereology}

Technological systems introduce an engineered parthood regime. In artificial wholes, parts are defined by interface compatibility rather than natural constraint alone.

\begin{definition}[Interface-bound whole]
A whole $w$ is interface-bound if its admissible parts are determined by design constraints that specify replaceability conditions.
\end{definition}

This regime differs from biological and chemical regimes in that parthood may be deliberately reconfigured without destroying identity. In biological systems, reconfiguration is constrained by metabolic closure and viability conditions; in chemical systems, it is constrained by bonding structure and energy minima. By contrast, in technological systems, identity is often defined at the level of functional specification rather than material continuity. Components may be replaced, upgraded, or rearranged provided that interface conditions and performance constraints are preserved.

The MSO therefore distinguishes between natural closure and engineered closure. Natural closure refers to self-maintaining constraint systems whose persistence depends upon internally regulated flows and boundary conditions. Such systems maintain their own coherence through dynamical feedback, and their identity is inseparable from their metabolic or thermodynamic organization. Engineered closure, by contrast, refers to design-maintained constraint systems in which coherence is externally specified through interface standards, modular architectures, and replaceability criteria. In this regime, the persistence of the whole is secured not by spontaneous self-maintenance alone but by adherence to an abstract specification that governs admissible reconfiguration.

The presence of technological glyphs within the MSO underscores that artificial wholes are not merely derivative or metaphorical entities. They constitute a distinct ontological stratum within the scale hierarchy, characterized by composition under specification and persistence under controlled substitution. Their inclusion in the MSO affirms that engineered systems participate fully in the layered ontology, even though the principles of their closure differ structurally from those governing natural regimes.

\section{Constellational Partitioning and Cognitive Imposition}

The constellational glyphs signal that some wholes arise from pattern-recognition imposed by observers. Constellations are epistemic partitions projected onto a continuous stellar distribution.

\begin{definition}[Imposed whole]
A whole is imposed if its unity derives from cognitive or symbolic stabilization rather than from a dynamical binding relation.
\end{definition}

The MSO does not treat imposed wholes as illusions. Rather, it treats them as legitimate strata in the epistemic ordering, provided their status is made explicit.

\section{Entropy and Stability Across Scales}

A final structural principle underlying the MSO is stability under perturbation.

\begin{definition}[Stable whole]
A whole $w$ at scale $s$ is stable if small perturbations to its parts do not destroy the invariants that define $w$ at that scale.
\end{definition}

This criterion permits cross-scale comparison. Atomic identity is stable under many chemical perturbations; biological identity is stable under cellular turnover; planetary identity is stable under seasonal variation.

Thus the scale-indexed ladder can be partially ordered by stability class rather than by mere size.

\section{The MSO as Bidirectional Ladder}

The Mereological Space Ontology proposes that the universe admits a stratified description in which part--whole relations are indexed by scale and stabilized by constraint. The glyph-string is not a decorative narrative but a compact encoding of regime transitions.

Read forward, it articulates an ascent through compositional stabilization, culminating in horizon and boundary structure.

Read backward, it narrates epistemic descent from cosmographic partitioning through planetary and biological regimes to microphysical mediation and finally to the observer.

The philosophical thesis is not that one direction is fundamental. It is that both are consistent only if one distinguishes carefully between ontic structure and observer-induced structure, and if one permits regime-boundary transitions in which new invariants define admissible wholes.

The MSO is therefore neither strict reductionism nor unstructured holism. It is a disciplined account of layered composition under constraint.

\section{Connection to the Relativistic Scalar--Vector Plenum (RSVP)}

The Mereological Space Ontology (MSO) may be understood as a discrete, scale-indexed articulation of a deeper continuous substrate. Within the RSVP framework, the universe is modeled not as a collection of primitive objects but as a structured plenum described by coupled scalar and vector fields together with entropy flow. The MSO can therefore be interpreted as the stratified stabilization of field configurations into admissible part--whole regimes.

Let $(\Phi, \mathbf{v}, S)$ denote the scalar density, vector flow, and entropy fields of RSVP. In the plenum description, these fields are fundamental; objects emerge as relatively stable regions in which field configurations satisfy constraint equations of the general form
\[
\mathcal{F}(\Phi, \mathbf{v}, S) = 0,
\]
together with boundary conditions preserving coherence under perturbation.

From this perspective, a ``whole'' in the MSO corresponds to a region $W \subset \mathbb{R}^4$ (or appropriate manifold) such that the restricted fields $(\Phi|_W, \mathbf{v}|_W, S|_W)$ exhibit dynamical closure. Closure here means that internal flows dominate over external flux sufficiently to preserve a recognizable invariant across time.

\subsection{Mereological Stabilization as Entropic Constraint}

In RSVP, structure arises through entropic differential and gradient regulation. A region becomes ontologically salient when entropy curvature and flow constraints produce a bounded instability interval followed by nonlinear saturation. Such regions are neither arbitrary partitions nor purely imposed descriptions; they are attractor configurations in field space.

Formally, suppose that for a region $W$ there exists a Lyapunov functional $\mathcal{L}[W]$ such that
\[
\frac{d}{dt} \mathcal{L}[W] \leq 0,
\]
with equality achieved only at equilibrium or steady-state flow. Then $W$ qualifies as a dynamically coherent unit. The MSO’s notion of a stable whole is thus grounded in RSVP as a region of entropically regulated persistence.

Atomic, molecular, biological, planetary, and galactic units correspond to different parameter regimes of the same underlying field equations. What the MSO encodes discretely as glyph-marked regime transitions corresponds in RSVP to bifurcations or effective-field approximations across scales.

\subsection{Scale-Indexed Mereology as Effective Field Theory}

The family $\{(U_s, \leq_s)\}_{s \in S}$ of scale-indexed mereologies can be mapped onto RSVP by treating each scale $s$ as an effective description obtained through coarse-graining of $(\Phi, \mathbf{v}, S)$. Let $\mathcal{C}_s$ denote a coarse-graining operator. Then
\[
(\Phi_s, \mathbf{v}_s, S_s) = \mathcal{C}_s(\Phi, \mathbf{v}, S).
\]

At each scale, admissible wholes are determined not by exhaustive microscopic specification but by invariants preserved under the coarse-graining operator $\mathcal{C}_s$. The passage from one regime to another therefore corresponds to a shift in which terms of the effective dynamics dominate and which quantities remain stable under projection. In this sense, regime boundaries in the MSO are not arbitrary thresholds but reflect qualitative reorganizations of the governing dynamics.

At microphysical scales, the effective description is dominated by local field interactions and exchange-mediated processes, so that admissible units are identified through symmetry structure and interaction vertices rather than through macroscopic persistence. At chemical scales, the salient invariants are bonding minima in scalar potential landscapes, and composition is governed by valence constraints and energy minimization principles. At biological scales, the defining feature of a whole is sustained entropy export and the maintenance of boundary conditions under continuous flow, so that organization is functional and thermodynamic rather than purely structural. At planetary scales, stabilization emerges from large-scale gradient balancing and energy throughput across layered subsystems, yielding coherent geophysical and ecological units. At cosmological scales, the relevant invariants shift again toward global curvature, large-scale flow alignment, and filamentary distribution of density, where wholes are defined by dynamical coherence across vast separations.

Under this interpretation, the MSO can be understood as a layered effective-field ontology extracted from the RSVP substrate. Each stratum corresponds to a regime in which particular terms in the coupled field dynamics become dominant, and in which a distinct class of invariants defines what qualifies as a persistent whole. The discrete ladder of the MSO thus reflects a sequence of effective descriptions grounded in a single continuous plenum.

\subsection{Observer as Field Subconfiguration}

Within RSVP, the observer is not external to the plenum. The glyph marking observation corresponds to a localized high-coherence configuration of $(\Phi, \mathbf{v}, S)$ capable of modeling its own environment. The epistemic ladder of the MSO therefore becomes an internal traversal within the same field substrate.

Observation induces an epistemic coarse-graining $\mathcal{O}$ that partitions the plenum according to the informational capacities of the observer-configuration. Ontic mereology and epistemic mereology coincide when observational partitions align with dynamically stable attractors.

Misalignment produces imposed wholes, symbolic partitions, or constellational structures that are epistemically useful but not dynamically fundamental.

\subsection{Filaments and Lamphrodynamic Connectivity}

The filamentary glyph in the MSO corresponds naturally to RSVP’s lamphrodynamic picture in which energy differentials propagate through structured flows without global metric expansion. Cosmic filaments are interpreted not merely as gravitational aggregations but as large-scale entropy-gradient channels.

In RSVP terms, these structures correspond to extended regions where $\mathbf{v}$ aligns with density gradients and entropy descent is coordinated across scales. Mereological fusion at cosmological scales is therefore a consequence of flow alignment rather than static aggregation.

\subsection{Horizon and Boundary Conditions}

The MSO glyph marking the threshold or horizon aligns with RSVP’s treatment of global boundary conditions. The observable universe is defined not by absolute edge but by causal and entropic constraints on information propagation.

In formal terms, let $\mathcal{H}$ denote a horizon region beyond which signal integration is dynamically suppressed. Then the MSO’s terminal boundary glyph corresponds to the maximal domain over which coherent integration of $(\Phi, \mathbf{v}, S)$ is possible for a given observer-configuration.

\subsection{Synthesis}

The structural relationship between the MSO and the RSVP framework may be articulated without appeal to schematic enumeration by emphasizing their complementary levels of description. The RSVP framework furnishes the continuous field-theoretic substrate, modeling the universe as a coupled scalar--vector--entropy plenum whose dynamics determine the evolution of density, flow, and thermodynamic constraint. Within this substrate, no discrete objects are assumed as primitive; instead, structure arises through regulated instability, saturation, and persistent flow alignment.

The MSO, by contrast, provides a discrete articulation of these stabilized regimes. It does not posit an independent ontology layered atop the plenum, but rather offers a scale-indexed description of those regions in field space that exhibit sufficient dynamical closure to qualify as wholes. In this sense, the MSO is a coarse-grained shadow of RSVP dynamics, organizing stabilized configurations into a stratified mereological hierarchy.

Stable wholes in the MSO correspond, under this interpretation, to entropically regulated attractors in the RSVP field dynamics. A whole is not defined by static containment but by persistence under flow, which in RSVP terms means that entropy production and gradient alignment are locally balanced so as to maintain structural invariants across time.

Regime transitions in the MSO correspond to effective-theoretic shifts in RSVP. When the dominant terms in the coupled field equations change, or when bifurcation phenomena reorganize attractor structure, the admissible units at the descriptive level must also change. What appears in the MSO as a boundary between scales is, in RSVP, a change in dynamical dominance or constraint architecture.

Finally, the observer-indexed dimension of the MSO corresponds in RSVP to internal coarse-graining within the plenum. The observer is itself a stabilized configuration of $(\Phi,\mathbf{v},S)$, and the epistemic partitioning it induces is a particular projection of the continuous field onto a representation space. Thus the ontic and epistemic layers are not external to one another but are related by internal transformations of the same underlying substrate.

Taken together, this synthesis shows that the MSO and RSVP are not rival ontologies but different resolutions of a single structured reality: RSVP supplies the continuous dynamical foundation, while the MSO organizes its stabilized configurations into a coherent mereological architecture indexed by scale.

The philosophical consequence is that the universe need not be conceived as a stack of independently fundamental layers. Instead, it is a single structured plenum whose dynamically stable regions admit layered description. The MSO is thus the mereological shadow cast by RSVP’s scalar--vector--entropy dynamics when viewed through scale-indexed coarse-graining.

In this way, the MSO and RSVP are not competing ontologies but complementary projections: one discrete and stratified, the other continuous and field-theoretic, each clarifying the other’s structural commitments.

\section{Measure-Theoretic Stability Criterion}

The MSO notion of a \emph{stable whole} may be made precise by treating each scale regime as carrying a measurable structure that supports both coarse-graining and perturbation analysis. The point is to define stability in a way that is invariant under relabeling of microstates and robust under the coarse-graining maps that relate adjacent scales.

Fix a scale $s \in S$. Let $(X_s,\Sigma_s,\mu_s)$ be a measure space of states at scale $s$, where $X_s$ is a configuration space, $\Sigma_s$ a $\sigma$-algebra of measurable events, and $\mu_s$ a reference measure encoding typicality (or equilibrium, or prior weight) at that scale. Let $\mathcal{C}_s : X \to X_s$ be a coarse-graining from an underlying fine-scale state space $X$ (when a common fine space is assumed), and let $\mathbb{P}_s$ be a Markov semigroup (or more generally a measurable evolution operator) acting on probability measures on $X_s$.

A candidate \emph{whole} at scale $s$ will be represented by a measurable subset $W_s \in \Sigma_s$, interpreted as the set of states that instantiate that whole. Stability then concerns the persistence of probability mass in $W_s$ under evolution and under small perturbations of initial conditions.

\begin{definition}[Measure-theoretic persistence]
Given $\epsilon \in (0,1)$ and a time horizon $T>0$, a measurable set $W_s \in \Sigma_s$ is \emph{$(\epsilon,T)$-persistent} if for every initial probability measure $\nu$ with $\nu(W_s)\ge 1-\epsilon$, one has
\[
(\nu \,\mathbb{P}_s^{t})(W_s) \ge 1-\epsilon
\quad \text{for all } t \in [0,T].
\]
\end{definition}

Persistence captures the idea that if the system is ``in'' the whole with high probability, it remains in that whole with high probability over the relevant time scale.

To distinguish true ontological stability from brittle classification, one must also demand robustness under perturbations. A convenient perturbation notion is Wasserstein distance (when $X_s$ has a metric) or total variation distance (in purely measure-theoretic settings).

\begin{definition}[Robust stability under perturbation]
Assume $X_s$ is a Polish space and let $W_1$ be the $1$-Wasserstein distance. A measurable set $W_s$ is \emph{$\delta$-robustly $(\epsilon,T)$-persistent} if there exists $\delta>0$ such that for every pair of initial measures $\nu,\tilde{\nu}$ with
\[
W_1(\nu,\tilde{\nu}) \le \delta
\quad \text{and} \quad
\nu(W_s)\ge 1-\epsilon,
\]
one has
\[
(\tilde{\nu}\,\mathbb{P}_s^{t})(W_s)\ge 1-\epsilon
\quad \text{for all } t \in [0,T].
\]
\end{definition}

This definition formalizes a strong form of ``small changes in parts do not destroy the whole,'' now expressed as stability of mass under evolved distributions.

The MSO can interpret parthood as induced by stability classes: if $A_s,B_s\in\Sigma_s$ represent candidate units, one can define an \emph{effective parthood} relation by requiring that typical realizations of $A_s$ lie within realizations of $B_s$ up to small measure error. For instance, one may set
\[
A_s \preccurlyeq_{\eta} B_s
\quad \Longleftrightarrow \quad
\mu_s(A_s \setminus B_s) \le \eta\,\mu_s(A_s),
\]
for a tolerance $\eta \in [0,1)$. This recovers the idea that at a given scale, units are defined by stable invariants rather than exact set inclusion.

\begin{remark}
The measure $\mu_s$ should be understood as regime-relative: in a thermodynamic regime it might be a Gibbs measure; in a dynamical regime it might be an SRB measure; in an epistemic regime it might be a prior induced by instrumentation. The MSO does not require a unique canonical choice, but it does require that stability be stated relative to an explicit measure.
\end{remark}

\section{Spectral and Topological Treatment of Regime Transitions}

The MSO treats regime transitions as changes in admissible units and invariants. Two complementary mathematical viewpoints clarify this: a spectral view in which transitions correspond to changes in dominant modes of evolution, and a topological view in which transitions correspond to changes in the shape (homotopy type) of attractor sets or critical loci.

\subsection{Spectral View: Mode Selection and Gap Opening}

Let $\mathcal{L}_s$ be a linearized generator of dynamics at scale $s$ acting on a suitable function space $\mathcal{H}_s$ (e.g.\ $L^2(X_s,\mu_s)$). One may think of $\mathcal{L}_s$ as the infinitesimal generator of $\mathbb{P}_s^t$. The spectrum $\sigma(\mathcal{L}_s)$ governs growth, decay, and oscillatory response.

A regime is characterized by a class of observables $\mathcal{A}_s \subset \mathcal{H}_s$ whose evolution is effectively governed by a small number of dominant spectral components of the generator $\mathcal{L}_s$. In such a situation, the long-time behavior of the system can be approximated by projection onto a finite-dimensional subspace spanned by eigenfunctions or generalized eigenmodes corresponding to isolated portions of the spectrum. The regime is therefore defined not merely by scale, but by the existence of a coherent spectral hierarchy separating relevant from irrelevant degrees of freedom.

A regime transition occurs when this spectral organization changes qualitatively. One paradigmatic case is a sign change in the real part of a leading eigenvalue, signaling the onset of instability and the breakdown of the previous stability structure. Another is a modification of the spectral gap, namely a change in the separation between slow and fast modes, which alters the effective dimensionality of the reduced dynamics and thereby shifts the appropriate ontology of persistent units. A transition may also occur through migration of spectral weight between families of modes, indicating that previously subdominant degrees of freedom become dynamically significant, thereby redefining which observables qualify as regime-defining. Finally, the emergence of continuous spectrum, or the dissolution of discrete eigenvalue structure, signals a loss of simple modal decomposition and may correspond to the disappearance of sharply defined wholes at that scale.


\begin{definition}[Spectral regime transition]
A transition $s\to t$ is spectral if there exists a parameter $\lambda$ controlling coarse-graining or effective description such that at $\lambda=\lambda_\ast$ the spectrum of $\mathcal{L}_\lambda$ undergoes a qualitative change, e.g.\ a leading eigenvalue crosses the imaginary axis or a spectral gap collapses.
\end{definition}

In each of these cases, the spectral data reorganize in such a way that the previous reduction to a small set of dominant modes is no longer valid. The ontology of stable units must then be reformulated in terms of the new spectral hierarchy, marking a genuine regime boundary within the MSO framework.

In MSO terms, this is precisely the moment at which ``what counts as a unit'' changes because the stability timescales and relevant invariants change. A new regime is one in which a different set of slow modes becomes dominant, permitting a new ontology of persistent wholes.

\subsection{Topological View: Attractors, Stratification, and Critical Loci}

The spectral view describes stability in terms of linear modes. The topological view describes it in terms of the geometry of invariant sets.

Let $\Phi_s^t : X_s \to X_s$ be the flow (or evolution map) at scale $s$ and let $\mathcal{A}_s \subset X_s$ be an attractor or invariant set supporting typical behavior (with respect to $\mu_s$). A regime transition can be interpreted as a change in the topological type of $\mathcal{A}_s$ or of a critical locus in an associated variational principle.

\begin{definition}[Topological regime transition]
A transition is topological if there exists a parameter $\lambda$ such that the invariant set $\mathcal{A}_\lambda$ changes its topological invariants at $\lambda=\lambda_\ast$, for example its number of connected components, its Betti numbers, or its homotopy type.
\end{definition}

A convenient sufficient condition is a bifurcation that changes the number or stability of equilibria. If the dynamics arise from (or admit) a Lyapunov functional $E_\lambda : X_s \to \mathbb{R}$, then the critical locus
\[
\mathrm{Crit}(E_\lambda) = \{x \in X_s : \nabla E_\lambda(x)=0\}
\]
may itself change topology as $\lambda$ crosses a threshold. In Morse-theoretic language, changes in the index and multiplicity of critical points produce changes in sublevel-set topology. In MSO terms, this corresponds to the emergence or disappearance of admissible ``wholes'' as stable basins in state space.

\subsection{Stratified State Spaces and Regime Sheaves}

Many MSO transitions are best modeled not by a single smooth manifold of states but by a \emph{stratified} space: different regimes correspond to different strata glued along singular boundaries. Let $X$ be stratified as
\[
X = \bigsqcup_{\alpha \in I} X_\alpha,
\]
with each stratum $X_\alpha$ carrying its own effective dynamics and invariants. A regime transition is then the passage of typical trajectories from one stratum to another, often through a singular set where the effective description changes.

This viewpoint naturally matches the MSO claim that the ontology is layered: each scale regime is a stratum with its own admissible unit language, and the ``glyph'' marking a regime boundary can be read as the symbol of a stratum change.

\subsection{Synthesis: Two Criteria for MSO Regime Boundaries}

The MSO can therefore operationalize a regime boundary between adjacent scales $s$ and $t$ by appealing to distinct but complementary mathematical diagnostics. One such diagnostic is spectral: a regime transition is indicated when the linearized generator acquires a new set of slow modes, or when an existing spectral gap collapses, thereby altering the hierarchy of dominant timescales and forcing a redefinition of the effective unit language. In such a case, previously negligible degrees of freedom become dynamically relevant, and the ontology of stable wholes must be reorganized accordingly. A second diagnostic is topological: a regime transition is signaled when the invariant set supporting typical behavior undergoes a qualitative change in topology, for example through bifurcation, the appearance or disappearance of connected components, or a shift in homotopy type. Such a change reflects a reconfiguration in the space of stable configurations itself, marking a structural boundary in the stratified ontology. Together, these diagnostics provide mathematically precise criteria for identifying genuine regime transitions within the MSO framework.


These diagnostics are compatible with the earlier measure-theoretic stability criterion. In practice, one expects that when a regime transition occurs, persistent sets $W_s$ fail to remain persistent under the previous regime’s evolution operator, and a new family of persistent sets emerges under the post-transition dynamics.

The philosophical consequence is that the MSO does not treat the scale ladder as an arbitrary hierarchy of ``smaller to bigger,'' but as a sequence of stability regimes whose transitions are mathematically detectable either spectrally or topologically.

\section{A Category-Theoretic Reformulation of the MSO}

The Mereological Space Ontology (MSO) may be reformulated categorically by replacing the primitive parthood relation with structured morphisms and by treating scale as an indexing parameter in a fibered categorical system. This recasts the MSO as a layered compositional architecture in which regime transitions correspond to changes in categorical structure rather than merely to changes in object size.

\subsection{Mereology as a Thin Category}

A classical mereological structure $(U,\leq)$ may be regarded as a thin category $\mathcal{M}$ in which the objects are the entities of $U$ and in which there exists a unique morphism $x \to y$ precisely when $x \leq y$. In this formulation, composition of morphisms corresponds directly to the transitivity of parthood: if $x \leq y$ and $y \leq z$, then the composite morphism $x \to z$ expresses $x \leq z$. Identity morphisms correspond to reflexivity, since each entity is a part of itself.

Because $\mathcal{M}$ is thin, there is at most one morphism between any two objects, and the categorical structure is completely determined by the partial order. In this way, classical mereology is recovered as a special case of category theory in which inclusion relations are reinterpreted as arrows and transitivity is internalized as composition.

Fusion, when it exists, admits a categorical interpretation as a colimit. More specifically, given a diagram of parts whose inclusion relations are encoded in $\mathcal{M}$, a fusion of those parts corresponds to a colimit object satisfying the universal property that every cocone from the diagram factors uniquely through it. In the posetal case, such a colimit reduces to a least upper bound or join. Thus the notion of a whole as the minimal entity containing given parts is precisely captured by the universal characterization of a colimit in a thin category.

\begin{definition}[Fusion as colimit]
Given a diagram $D : J \to \mathcal{M}$ of parts, a fusion $\mathrm{colim}(D)$, when it exists, is a colimit in $\mathcal{M}$ satisfying the universal property that every cocone from $D$ factors uniquely through it.
\end{definition}

Thus classical mereology corresponds to a category enriched in truth values (a preorder), and wholes correspond to universal constructions.

\subsection{Scale as an Indexing Category}

Let $\mathcal{S}$ be a totally ordered category of scales. The MSO may then be represented as a functor
\[
\mathcal{F} : \mathcal{S}^{op} \to \mathbf{Cat},
\]
assigning to each scale $s$ a category $\mathcal{M}_s$ of entities at that scale. For $s \leq t$, the morphism $s \to t$ in $\mathcal{S}$ induces a functor
\[
\mathcal{F}(t \to s) = \pi_{t\to s} : \mathcal{M}_t \to \mathcal{M}_s,
\]
interpreted as coarse-graining.

This construction makes the MSO a contravariant diagram of categories indexed by scale. Ontic ascent corresponds to moving forward in $\mathcal{S}$, while epistemic descent corresponds to applying the functors $\pi_{t\to s}$.

\subsection{Regime Transitions as Change of Monoidal Structure}

Many MSO strata admit not merely a partial order but a symmetric monoidal structure $(\mathcal{M}_s,\otimes_s,\mathbf{1}_s)$ representing composition of parts.

A regime transition $s\to t$ may involve a change in the monoidal product:

\[
(\mathcal{M}_s,\otimes_s) \quad \not\simeq \quad (\mathcal{M}_t,\otimes_t).
\]

At microphysical scales, the monoidal product $\otimes_s$ may be interpreted as field-theoretic interaction, in which composition corresponds to the coupling of local degrees of freedom under exchange-mediated dynamics. In such a regime, fusion is governed by symmetry constraints and conservation laws rather than by macroscopic adjacency.

At chemical scales, the same formal symbol $\otimes_s$ acquires a different operational meaning: it denotes bonding operations constrained by valence structure and potential minima in an effective scalar landscape. Composition here is neither arbitrary aggregation nor mere juxtaposition, but the stabilization of electron-sharing or exchange configurations that generate discrete molecular units.

At biological scales, $\otimes_s$ must be read as functional integration under sustained entropy export. Composition is defined not solely by structural bonding but by participation in metabolic closure and regulatory coupling. A biological whole persists insofar as its parts co-maintain boundary conditions against environmental flux.

At technological scales, the monoidal product $\otimes_s$ becomes interface composition, in which admissible fusion is determined by design constraints and compatibility conditions. The identity of the whole is preserved under replacement of parts provided that interface invariants are respected. In this regime, composition is mediated by specification and protocol rather than by purely natural constraint.

\begin{definition}[Monoidal regime transition]
A regime transition is monoidal if there is no monoidal equivalence between $(\mathcal{M}_s,\otimes_s)$ and $(\mathcal{M}_t,\otimes_t)$ preserving relevant invariants.
\end{definition}

Thus regime boundaries correspond categorically to structural changes in how composition is defined.

\subsection{Stability as a Subcategory of Persistent Objects}

Let $\Phi_s$ denote the evolution functor (or endofunctor) at scale $s$. Stability can be expressed categorically by defining a full subcategory
\[
\mathrm{Stab}(\mathcal{M}_s) \subseteq \mathcal{M}_s
\]
consisting of objects $W$ such that $\Phi_s(W)$ is isomorphic to $W$ up to specified tolerance or equivalence.

\begin{definition}[Persistent object]
An object $W \in \mathcal{M}_s$ is persistent if there exists a natural isomorphism
\[
\eta_W : \Phi_s(W) \to W.
\]
\end{definition}

The MSO’s stable wholes are precisely these persistent objects. Regime transitions correspond to changes in the subcategory $\mathrm{Stab}(\mathcal{M}_s)$.

\subsection{Filamentary Connectivity as Higher Morphisms}

Large-scale connectivity and the existence of non-local wholes motivate extending the categorical formulation of the MSO beyond thin categories to higher-categorical structure. In a purely thin category, parthood is encoded as a single layer of morphisms, and there is no intrinsic representation of the relations that bind otherwise distinct wholes into coherent systems. However, filamentary and network-like regimes indicate that relational bindings themselves possess structure and coherence conditions that cannot be reduced to simple inclusion.

In a 2-category formulation, objects correspond to entities or stabilized wholes at a given scale. 1-morphisms encode parthood or inclusion relations, preserving the classical mereological interpretation at the base level. In addition to these, 2-morphisms represent structured relational bindings between 1-morphisms. Such 2-morphisms may correspond, depending on regime, to gravitational coupling, informational transfer, functional coordination, or any constraint that binds two parts into a higher-order coherence without necessarily fusing them into a single contiguous object.

Under this interpretation, a filament is not merely a composite object but a coherent family of 1-morphisms equipped with 2-morphisms satisfying compatibility conditions. The coherence axioms of the 2-category then express the fact that relational binding is not arbitrary: compositions of bindings must associate and compose consistently. This captures the MSO intuition that certain large-scale wholes are stabilized not by adjacency but by structured interaction.

Formally, let $\mathcal{M}_s$ be the category of entities at scale $s$, and consider a 2-category $\mathbf{M}_s$ whose objects are those of $\mathcal{M}_s$, whose 1-morphisms are parthood maps, and whose 2-morphisms encode relational constraints between such maps. The presence of nontrivial 2-morphisms allows the representation of networks and webs as higher-order coherence structures rather than as mere unions. A cosmic filament, for example, may be represented as a diagram of objects and parthood maps together with specified 2-morphisms expressing dynamical alignment or constraint propagation.

Filamentary structures correspond to coherent systems of 2-morphisms binding otherwise separated objects. The cosmic web thus becomes a higher-categorical object: not merely a fusion but a network of relations with coherence conditions.

In this higher-categorical setting, regime transitions correspond not only to changes in objects and monoidal structure but also to changes in the permissible 2-morphisms and their coherence laws. A scale boundary may therefore be characterized by a shift in which relational bindings are admissible or dominant. The MSO thus generalizes from a mereology of inclusion to a stratified higher-category in which relational structure itself becomes a first-class ontological component.

\subsection{Observer as Internal Functor}

Within this categorical framework, the observer corresponds to an internal functor
\[
\mathcal{O}_s : \mathcal{M}_s \to \mathbf{Rep},
\]
where $\mathbf{Rep}$ is a category of representations (states, signals, measurements). Epistemic mereology is induced by factorization through $\mathcal{O}_s$.

Ontic structure corresponds to morphisms in $\mathcal{M}_s$; epistemic structure corresponds to morphisms preserved or reflected by $\mathcal{O}_s$.

A whole may be ontically real but epistemically invisible if $\mathcal{O}_s$ fails to preserve its defining morphisms.

\subsection{Regime Boundaries as Changes in Fibered Structure}

Combining the above, the MSO may be viewed as a category fibered over $\mathcal{S}$:

\[
p : \mathcal{E} \to \mathcal{S},
\]

where the fiber over $s$ is $\mathcal{M}_s$. A regime transition corresponds to a change in the fiber structure that is not trivializable by equivalence.

If $p$ fails to admit a cartesian lifting preserving monoidal structure across $s\to t$, the transition marks a genuine ontological boundary.

\subsection{Synthesis}

The category-theoretic reformulation of the MSO permits a concise structural interpretation in which the traditional vocabulary of parts and wholes is recast in categorical language. Parthood is no longer treated as a primitive binary relation but as a morphism within an appropriate category, so that inclusion is represented by arrow structure rather than by set-theoretic containment. Fusion, when it exists, is expressed as a colimit, thereby capturing the universal property that characterizes a whole as the minimal object through which all constituent inclusions factor.

Scale is represented by an indexing category whose objects correspond to regimes of description and whose morphisms encode admissible transitions between them. The MSO then appears as a diagram of categories fibered over this indexing structure. Coarse-graining is realized as a functor between fibers, typically contravariant with respect to the scale ordering, and it preserves only those invariants appropriate to the target regime.

Stability is interpreted as persistence under an endofunctor describing temporal or dynamical evolution. Objects that admit a natural isomorphism to their image under this endofunctor qualify as persistent wholes at the given scale. Regime transitions correspond to structural changes within fibers, including alterations in monoidal structure, spectral organization of generators, or topological configuration of invariant sets. Finally, the observer is modeled as an internal representation functor from a regime category to a category of representations, thereby inducing an epistemic mereology derived from the preservation and reflection of morphisms.

Taken together, these identifications present the MSO as a stratified categorical architecture in which compositional laws, stability conditions, and observational structure are unified within a single formal system indexed by scale.

In this formulation, the MSO is not merely a hierarchical ordering of objects but a stratified categorical system in which compositional laws themselves evolve across scales. The glyph-string becomes a compact encoding of transitions between categorical fibers, each with its own admissible morphisms, invariants, and universal constructions.

Thus the MSO may be interpreted as a fibered, monoidal, scale-indexed category of structured wholes emerging from a continuous substrate and stabilized by constraint.

\section{An Interpretation of the MSO as a Finite Lie Algebra}

The MSO glyph-string may be interpreted not only as an ordered presentation of regimes but also as a finite algebra of regime-operators whose noncommutativity encodes the empirical fact that ``changing scale'' and ``changing descriptive primitive'' do not, in general, commute. A compact way to formalize this noncommutativity is to regard the glyphs as a basis of a finite-dimensional real Lie algebra whose bracket measures obstruction to simultaneous stabilization of two regime moves.

\subsection{The MSO Lie Algebra of Regime Operators}

Let $\Sigma$ be the finite set of glyphs used by the MSO chain and write $\mathbb{R}\Sigma$ for the real vector space spanned by formal basis elements $\{e_g : g \in \Sigma\}$. Define a Lie algebra structure on $\mathfrak{g}_{\mathrm{MSO}} := (\mathbb{R}\Sigma,[\cdot,\cdot])$ by specifying a bilinear antisymmetric bracket satisfying Jacobi. The intended reading is that $e_g$ represents the infinitesimal generator of a regime transformation $T_g$ acting on a space of descriptions (or observables, or coarse-grainings).

One may motivate this by imagining a local group action on a description space $\mathcal{D}$: for small parameters $\epsilon$, the operator $\exp(\epsilon e_g)$ acts as a regime perturbation. Then the commutator
\[
[e_g,e_h]
\]
measures the failure of applying $g$ then $h$ to be equivalent to applying $h$ then $g$. In MSO language, this bracket records the obstruction to making two regime cuts simultaneously stable.

\subsection{A Graded Structure by Scale}

The MSO ladder already carries an implicit scale ordering. This may be expressed as a $\mathbb{Z}$-grading on $\mathfrak{g}_{\mathrm{MSO}}$. Choose a degree map
\[
\deg : \Sigma \to \mathbb{Z}
\]
that increases monotonically along the forward ontic chain (or decreases, depending on convention). Extend $\deg$ linearly so that $\mathfrak{g}_{\mathrm{MSO}}$ decomposes as
\[
\mathfrak{g}_{\mathrm{MSO}} = \bigoplus_{k\in\mathbb{Z}} \mathfrak{g}_k,
\qquad
\mathfrak{g}_k := \mathrm{span}\{e_g : \deg(g)=k\}.
\]
A natural constraint compatible with the MSO intuition is that brackets respect the grading:
\[
[\mathfrak{g}_i,\mathfrak{g}_j] \subseteq \mathfrak{g}_{i+j}.
\]
This condition states that composing a degree-$i$ and degree-$j$ perturbation yields an obstruction of combined ``scale weight.'' It is a formal surrogate for the empirical observation that cross-scale interactions often generate intermediate effective degrees of freedom.

\subsection{Structure Constants as Regime-Transition Couplings}

Fix a basis $\{e_g\}_{g\in\Sigma}$. Then the Lie bracket is determined by structure constants $C_{gh}^{k}$ defined by
\[
[e_g,e_h] = \sum_{k\in\Sigma} C_{gh}^{k}\, e_k,
\qquad
C_{gh}^{k} = - C_{hg}^{k}.
\]
The constants $C_{gh}^{k}$ are not arbitrary in the MSO reading. They encode which regime combinations produce a third regime effect. For example, it is reasonable to model that ``instrumental refinement'' combined with ``classification by intervention'' produces an obstruction that lives in a mediating regime rather than at the extremes. Similarly, ``technological accretion'' combined with ``observer partition'' may generate an epistemic bias term, capturing that observation is itself conditioned by tool-mediated representation.

Rather than committing to a unique empirical calibration, the MSO Lie algebra is best treated as a \emph{presentation} whose constants are selected to reflect desired compatibility constraints. The Jacobi identity then becomes a nontrivial consistency condition on how regime obstructions compose.

\begin{proposition}[Jacobi as coherence of regime obstructions]
For any $g,h,\ell \in \Sigma$, the Jacobi identity
\[
[e_g,[e_h,e_\ell]] + [e_h,[e_\ell,e_g]] + [e_\ell,[e_g,e_h]] = 0
\]
expresses the requirement that third-order incompatibilities between regime operations be coherent, in the sense that no cyclic composition yields a net ``anomaly'' in the algebra of regime cuts.
\end{proposition}

In practical terms, Jacobi encodes that if one passes between three regimes in different orders, the induced mismatch must be expressible entirely in terms of the same finite basis of regime effects.

\subsection{A Minimal Solvable Model with a Scale Derivation}

A particularly transparent finite Lie algebra model arises by separating the MSO into a ``scale traversal'' generator and a nilpotent ``regime shift'' ideal.

Let $d$ denote the basis element corresponding to the scale traversal glyph (the MSO $\mathrm{scale}$ operator), and let $\mathfrak{n}$ be the span of the remaining regime generators. Impose
\[
[d, x] = \lambda_{\deg(x)}\, x
\quad \text{for all } x \in \mathfrak{n},
\]
where $\lambda_{\deg(x)} \in \mathbb{R}$ is a weight depending only on the assigned degree. This makes $\mathrm{ad}_d$ act diagonally on $\mathfrak{n}$, so $d$ functions as a grading derivation. Then declare $\mathfrak{n}$ to be nilpotent by specifying brackets only among ``nearby'' regime generators and requiring that repeated commutators eventually vanish. The resulting algebra
\[
\mathfrak{g}_{\mathrm{MSO}} \cong \mathbb{R}d \ltimes \mathfrak{n}
\]
is solvable, and it captures a core MSO intuition: scale traversal reorganizes which regime operators are relevant, while local regime compositions generate higher-order, but eventually terminating, obstructions.

\subsection{A Path-Algebra Bracket from the MSO Chain}

Because the MSO chain is a linear ordering of glyphs, one may define a Lie bracket using adjacency on the chain. Let $\Sigma=\{g_1,\dots,g_N\}$ in chain order. A minimal nontrivial choice is to specify brackets only for adjacent pairs:
\[
[e_{g_i}, e_{g_{i+1}}] := \alpha_i\, e_{g_{i+1}},
\qquad
[e_{g_{i+1}}, e_{g_i}] := -\alpha_i\, e_{g_{i+1}},
\]
and set all other brackets to zero, then adjust (if necessary) by adding a small number of additional terms to satisfy Jacobi. This produces an ``upper-triangular'' algebra in which commutators push forward along the chain. Conceptually, it models the idea that combining a regime with its immediate refinement tends to produce an effect still controlled by the refined regime.

More generally, one may allow a short-range bracket
\[
[e_{g_i}, e_{g_j}] = 0 \quad \text{whenever } |i-j|>m,
\]
for a small interaction range $m$, reflecting that distant scales are weakly coupled at the level of regime obstructions except through chains of intermediate transitions.

\subsection{Interpretive Payoff}

Interpreting the MSO as a finite Lie algebra yields two structural consequences.

First, it turns the MSO into a noncommutative calculus of regime moves. The emphasis shifts from static classification to an algebra of transformations, consistent with the MSO theme that the world is carved by operators (coarse-graining, instrumentation, constraint selection) rather than by a fixed list of objects.

Second, it provides a clean interface to RSVP-style field ontology. In an RSVP setting, different effective regimes correspond to different dominant terms in a coupled $(\Phi,\mathbf{v},S)$ dynamics; the Lie bracket then represents the obstruction to simultaneously diagonalizing these effective descriptions. In that reading, $\mathfrak{g}_{\mathrm{MSO}}$ is a finite approximation to an infinite-dimensional symmetry algebra of coarse-graining and regime selection, truncated to the glyph basis.

The finite Lie algebra picture therefore complements the earlier measure-theoretic and spectral viewpoints: stability becomes the persistence of certain orbits of the induced action, regime transitions correspond to changes in the effective representation theory of $\mathfrak{g}_{\mathrm{MSO}}$, and the glyph-string is a distinguished word whose successive generators describe a path through the space of admissible decompositions.

\section{Conclusion: The MSO as a Discipline of Layered Description}

The Mereological Space Ontology has been developed as a compact but structurally rich account of how a universe may be organized into admissible parts and wholes across scales without collapsing into either naïve reductionism or unstructured holism. At its core lies a simple claim: parthood is always indexed by regime, and regime is always indexed by constraint and stability.

The glyph-string that animates the MSO is not a decorative taxonomy but a compressed encoding of regime operators. Read in the forward direction, it expresses an ontic ascent through progressively stabilized configurations, from observer-marked partitions through energetic mediation, atomic individuation, biological closure, planetary integration, technological accretion, and finally large-scale filamentary and horizon structure. Read in reverse, it expresses an epistemic descent through observational practice, instrumentation, classification, and theoretical refinement. The same ordered sequence supports both traversals because it is grounded in a single scale-indexed mereological scaffold.

Formally, this scaffold has been articulated in several complementary ways. The measure-theoretic treatment defines wholes as sets that are persistent and robust under regime-relative evolution operators. The spectral analysis characterizes regime boundaries as qualitative changes in dominant modes, spectral gaps, or the structure of generators. The topological analysis interprets transitions as bifurcations or changes in the invariant sets supporting typical behavior. The categorical reformulation recasts parthood as morphism, fusion as colimit, scale as an indexing category, and stability as persistence under an endofunctor. The Lie-algebraic interpretation treats regime operators as noncommuting generators whose brackets encode the obstructions and compatibilities of cross-scale transformations.

These formal lenses converge on a common structural thesis: regime transitions are not merely quantitative enlargements or refinements, but qualitative reorganizations of admissible invariants. A stable whole at one scale is not guaranteed to be definable in the same primitive language at another. Instead, each stratum admits its own closure conditions, monoidal composition laws, and persistence criteria. The MSO therefore insists that ontological articulation must be sensitive to the effective structure of each regime.

The connection to the Relativistic Scalar–Vector Plenum further clarifies the ontological modesty of the MSO. RSVP provides a continuous field-theoretic substrate in which stabilized regions of flow, density, and entropy appear as attractors. The MSO may then be understood as the discrete, scale-indexed shadow of these stabilized configurations under coarse-graining. Wholes are not metaphysically primitive; they are dynamically persistent subconfigurations. The layered ontology is not a stack of independently fundamental realms, but a hierarchy of effective descriptions extracted from a single structured plenum.

In this sense, the MSO functions as a discipline of description. It does not legislate what exists; it specifies the conditions under which something may count as a stable whole at a given scale. It distinguishes ontic mereology from epistemic mereology, natural closure from engineered closure, local adjacency from filamentary connectivity, and dynamical binding from imposed symbolic partition. It provides criteria—measure-theoretic, spectral, topological, categorical, algebraic—for recognizing genuine regime boundaries rather than merely apparent ones.

The philosophical consequence is a tempered realism about parts and wholes. Wholes are real insofar as they are stable under the appropriate regime of invariants; they are relative insofar as those invariants are scale-indexed and sometimes observer-mediated. The universe, under this view, is neither an undifferentiated plenum nor a mosaic of irreducible blocks. It is a single structured field whose stable configurations admit layered articulation.

The Mereological Space Ontology thus offers a unified account of composition, scale, and observation. Its glyph-string encodes a bidirectional ladder; its formal apparatus supplies criteria of stability and transition; its philosophical stance preserves both structural grounding and regime plurality. What emerges is not merely a taxonomy of levels, but a coherent architecture of part–whole relations in a layered, constraint-governed universe.

\newpage
\section*{Appendix} 

\appendix

\section*{MSO Symbol Table}

\begin{center}
\begin{tabular}{ll}
\textbf{Regime} & \textbf{Symbol} \\
\hline
ocularo & $\MSOglyph{Ocularo}$ \\
monoturno & $\MSOglyph{Monoturno}$ \\
lamphron & $\MSOglyph{Lamphron}$ \\
lamphrodyno & $\MSOglyph{Lamphrodyno}$ \\
gravito & $\MSOglyph{Gravito}$ \\
quarko & $\MSOglyph{Quarko}$ \\
electro & $\MSOglyph{Electro}$ \\
lepto & $\MSOglyph{Lepto}$ \\
fermio & $\MSOglyph{Fermio}$ \\
photo & $\MSOglyph{Photo}$ \\
gluo & $\MSOglyph{Gluo}$ \\
meso & $\MSOglyph{Meso}$ \\
boso & $\MSOglyph{Boso}$ \\
proto & $\MSOglyph{Proto}$ \\
neutro & $\MSOglyph{Neutro}$ \\
baryo & $\MSOglyph{Baryo}$ \\
hadro & $\MSOglyph{Hadro}$ \\
ato & $\MSOglyph{Ato}$ \\
moleculo & $\MSOglyph{Moleculo}$ \\
unicellulo & $\MSOglyph{Unicellulo}$ \\
crystolattisso & $\MSOglyph{Crystolattisso}$ \\
micro & $\MSOglyph{Micro}$ \\
megalobio & $\MSOglyph{Megalobio}$ \\
gaiama & $\MSOglyph{Gaiama}$ \\
technio & $\MSOglyph{Technio}$ \\
tella & $\MSOglyph{Tella}$ \\
gaiaselena & $\MSOglyph{Gaiaselena}$ \\
solara & $\MSOglyph{Solara}$ \\
perinebula & $\MSOglyph{Perinebula}$ \\
vacuo & $\MSOglyph{Vacuo}$ \\
oriocygnobrachio & $\MSOglyph{Oriocygnobrachio}$ \\
galacto & $\MSOglyph{Galacto}$ \\
proximasystada & $\MSOglyph{Proximasystada}$ \\
virgo & $\MSOglyph{Virgo}$ \\
laniakeasuperclusto & $\MSOglyph{Laniakeasuperclusto}$ \\
piscescetusfilamento & $\MSOglyph{Piscescetusfilamento}$ \\
cytosis & $\MSOglyph{Cytosis}$ \\
\end{tabular}
\end{center}

\newpage
\begin{thebibliography}{99}

\bibitem{Simons1987}
Peter Simons.
\newblock \emph{Parts: A Study in Ontology}.
\newblock Oxford University Press, 1987.

\bibitem{Varzi2016}
Achille C. Varzi.
\newblock Mereology.
\newblock In Edward N. Zalta (ed.), \emph{The Stanford Encyclopedia of Philosophy}, Winter 2016 Edition.

\bibitem{Lewis1991}
David Lewis.
\newblock \emph{Parts of Classes}.
\newblock Blackwell, 1991.

\bibitem{Fine1999}
Kit Fine.
\newblock Things and Their Parts.
\newblock \emph{Midwest Studies in Philosophy}, 23(1):61--74, 1999.

\bibitem{Chalmers2012}
David J. Chalmers.
\newblock \emph{Constructing the World}.
\newblock Oxford University Press, 2012.

\bibitem{LadymanRoss2007}
James Ladyman and Don Ross.
\newblock \emph{Every Thing Must Go: Metaphysics Naturalized}.
\newblock Oxford University Press, 2007.

\bibitem{Maudlin2007}
Tim Maudlin.
\newblock \emph{The Metaphysics Within Physics}.
\newblock Oxford University Press, 2007.

\bibitem{Tegmark2014}
Max Tegmark.
\newblock \emph{Our Mathematical Universe}.
\newblock Alfred A. Knopf, 2014.

\bibitem{Weinberg1995}
Steven Weinberg.
\newblock \emph{The Quantum Theory of Fields, Volume I}.
\newblock Cambridge University Press, 1995.

\bibitem{Carroll2010}
Sean M. Carroll.
\newblock \emph{From Eternity to Here: The Quest for the Ultimate Theory of Time}.
\newblock Dutton, 2010.

\bibitem{Prigogine1997}
Ilya Prigogine.
\newblock \emph{The End of Certainty: Time, Chaos, and the New Laws of Nature}.
\newblock Free Press, 1997.

\bibitem{Kauffman1993}
Stuart A. Kauffman.
\newblock \emph{The Origins of Order: Self-Organization and Selection in Evolution}.
\newblock Oxford University Press, 1993.

\bibitem{Barabasi2016}
Albert-László Barabási.
\newblock \emph{Network Science}.
\newblock Cambridge University Press, 2016.

\bibitem{Smolin2006}
Lee Smolin.
\newblock \emph{The Trouble with Physics}.
\newblock Houghton Mifflin, 2006.


\end{thebibliography}

\end{document}
